% Fault-tolerant migration protocol, organised around the replicated-log
% entries. Replaces design.tex:74-205 (kept there inside \iffalse ... \fi).
% \openq{...} marks decisions that are still open; it prints in red.
\providecommand{\openq}[1]{{\color{red}[OPEN: #1]}}

\subsection{Fault-tolerant migration protocol}
\label{sec:migration-protocol}

\paragraph{Roles and channels.}
A migration moves a set of slots from a \emph{donor} replica group to a \emph{recipient} replica group.
A coordinator $C$ splits the slots into \emph{migration groups} $G_1, \dots, G_n$ and instructs the donor leader to migrate them to the recipient leader; together, the groups form one migration transaction.
On each leader, the \emph{Migration Manager (MM)} drives the protocol.
The protocol uses three channels, which Fig.~\ref{fig:migration_flow} draws differently:
RPCs for small control messages,
one-sided RDMA writes for the bulk data,
and entries appended to each group's replicated (Raft) log for every decision that must survive a crash.
Bulk data never goes through the log.

\paragraph{Guarantees.}
Under the fault model of Section~\ref{sec:fault-model}, the protocol guarantees that
(i)~no committed write is lost,
(ii)~a read never observes partially migrated or unreplicated data, and
(iii)~a migration that has started rolls forward to completion; it is never rolled back.

\subsection{Migration log entries}
\label{sec:log-entries}

The progress of each group $G_i$ is recorded by five entries in the donor's and the recipient's Raft logs (Table~\ref{tab:log-entries}).
Every replica learns how far the migration got by applying its log, so a newly elected leader resumes from the last committed entry instead of relying on in-memory state, which a crash loses.

\begin{table}[t]
  \centering
  \footnotesize
  \caption{Replicated-log entries recording the migration of group $G_i$}
  \label{tab:log-entries}
  \setlength{\tabcolsep}{3pt}
  \begin{tabular}{@{}l l >{\raggedright\arraybackslash}p{2.3cm} >{\raggedright\arraybackslash}p{2.2cm}@{}}
    \hline
    \textbf{Entry} & \textbf{Log} & \textbf{Appended when} & \textbf{Enables} \\
    \hline
    \texttt{$G_i$\_START}           & donor     & before the transfer of $G_i$ & a new donor leader to resume $G_i$ \\
    \texttt{TXN\_START($G_i$)}      & recipient & before the recipient serves client operations on $G_i$ & tombstones for deleted keys \\
    \texttt{INDEX\_UPDATE($G_i$)}   & recipient & a majority holds the data and the leader's index update is done & double reads for $G_i$ to stop \\
    \texttt{TXN\_DONE($G_i$)}       & recipient & after \texttt{INDEX\_UPDATE} commits & a new recipient leader to skip $G_i$ \\
    \texttt{$G_i$\_DONE}            & donor     & the recipient reports $G_i$ complete & a new donor leader to skip $G_i$ \\
    \hline
  \end{tabular}
\end{table}

\paragraph{Durability invariant.}
The protocol rests on a single invariant:
\begin{quote}
  \itshape If \texttt{INDEX\_UPDATE($G_i$)} is committed, a majority of the recipient group holds the data of $G_i$.
\end{quote}
It holds because the recipient leader proposes \texttt{INDEX\_UPDATE($G_i$)} only after a majority of replicas have reported receiving the data (Section~\ref{sec:intra-recipient}).
The fault model guarantees that a majority of the group remains available; since any two majorities intersect, at least one available replica holds the data.
A newly elected leader that lacks part of it fetches it from the other replicas of its group before serving $G_i$.

\subsection{Migration phases}
\label{sec:migration-phases}

Fig.~\ref{fig:migration_flow} shows the migration of one group; the circled numbers below refer to its steps.
The transaction completes once every group has completed.

\subsubsection{Intra-donor preparation}
\label{sec:intra-donor}

\mstep{1}~The donor leader appends \texttt{$G_i$\_START} to its log and commits it, making the migration intent durable.
It then identifies the memory blocks holding the slots of $G_i$ and registers them with its RNIC as the source of the transfer.
\openq{item 2: when writes for $G_i$ start being redirected to the recipient, relative to \texttt{TXN\_START($G_i$)}.}

\subsubsection{Leader-to-leader transfer}
\label{sec:leader-to-leader}

\mstep{2}~The donor leader sends a \texttt{prepare} RPC to the recipient leader.
The recipient leader allocates landing buffers for $G_i$, registers them as memory regions, and replies with their addresses, sizes, and RDMA access keys (RKeys).
\mstep{3}~Using these keys, the donor leader writes the data of $G_i$ directly into the landing buffers with one-sided RDMA writes, which do not involve the recipient's CPU.
Once all writes complete, it sends an \texttt{RDMA\_DONE} RPC to the recipient leader.

\subsubsection{Intra-recipient replication}
\label{sec:intra-recipient}

\mstep{4}~On \texttt{RDMA\_DONE}, the recipient leader starts two activities in parallel: its own index update of $G_i$ (Section~\ref{sec:index_updating}), and the replication of the data to its followers.
Replication uses chain replication over RDMA.
The leader orders the live replicas into a chain and writes the data to the first one; each replica that has received the data forwards it to its successor and signals it with a \texttt{CHAIN\_DONE} message.
Replicas report which buffers they hold in their \texttt{AppendEntries} responses to the leader.

\mstep{5}~Once a majority of replicas, counting the leader, hold the data and the leader's index update has finished,
\mstep{6}~the leader proposes \texttt{INDEX\_UPDATE($G_i$)} through Raft.
Once it commits, the durability invariant holds for $G_i$, and the leader appends \texttt{TXN\_DONE($G_i$)} to close its side of the group.

\paragraph{Fast and slow paths.}
Replicating to a majority is the \emph{fast path}: it is synchronous and must finish before double reads for $G_i$ can stop.
The remaining replicas are on the \emph{slow path} and receive the data asynchronously, either from the ongoing chain or from the leader during log synchronization, when they apply an \texttt{AppendEntries} that contains \texttt{INDEX\_UPDATE($G_i$)}.
Eventually, all replicas converge to the same state.

\subsubsection{Index update}
\label{sec:index_updating}

Transferring the raw data is not enough: \OurImplementation{} keeps a dictionary-like index with pointers to the key-value pairs, which must be updated to point to the new memory locations.
This \emph{index update} runs on dedicated background threads with lower scheduling priority, to minimize interference with client-serving threads.
Since these threads run concurrently with client threads, \OurImplementation{} uses fine-grained locking for dictionary accesses.

\paragraph{Stale entries.}
During the migration, new writes for $G_i$ go to the recipient, so the transferred data may contain older versions of these keys.
The index update detects and discards such stale entries, keeping only the most recent value of each key.

\paragraph{Tombstones.}
A key $K$ that is updated and then deleted at the recipient before the index update runs would be resurrected from the donor's older copy.
To prevent this, recipient replicas keep tombstones for deleted keys while \texttt{TXN\_START($G_i$)} is active, and the index update skips any key with a tombstone.
The tombstones of $G_i$ are garbage-collected once its index update finishes.
\openq{item 3: confirm tombstones are implemented.}

\subsubsection{Completion}
\label{sec:commit-completion}

\mstep{7}~The recipient leader notifies the donor leader that $G_i$ is complete, and the donor leader appends and commits \texttt{$G_i$\_DONE}.
\openq{item 8: when ownership of the slots of $G_i$ is transferred to the recipient, per group or once at the end of the transaction.}

\subsection{Serving client requests during migration}
\label{sec:serving-during-migration}

\OurImplementation{} migrates data without interrupting client-facing operations.
Requests for slots that are not under migration are handled normally by their owning group.

\paragraph{Writes.}
Once the migration of $G_i$ begins, writes to its slots are redirected to the recipient group, processed by the recipient leader, and replicated through Raft.
The donor applies no further writes to $G_i$, so no write-back or reconciliation is needed after the migration; stale-entry detection and tombstones keep the migrated data from overwriting these writes.

\paragraph{Reads.}
Until \texttt{INDEX\_UPDATE($G_i$)} commits, clients use double reads, so that a read never depends on data the recipient may not yet hold.
Once it commits, the recipient group serves $G_i$ on its own: by the durability invariant, any newly elected leader can obtain the data from its group.

\subsection{Failure handling}
\label{sec:failure-handling}

A crash loses all migration state that is not in a log, such as landing buffers, chains, and in-flight RDMA writes.
Recovery therefore starts from the logs: a newly elected leader looks for groups whose start entry has no matching done entry and resumes them (Table~\ref{tab:failures}).
Recovery only moves forward, and resending the data of a group is harmless, because the index update discards entries that are not newer than what the recipient holds.

\begin{table}[t]
  \centering
  \footnotesize
  \caption{Recovery after a crash during the migration of $G_i$}
  \label{tab:failures}
  \setlength{\tabcolsep}{3pt}
  \begin{tabular}{@{}>{\raggedright\arraybackslash}p{1.45cm} >{\raggedright\arraybackslash}p{6.5cm}@{}}
    \hline
    \textbf{Crash} & \textbf{Recovery} \\
    \hline
    Donor follower     & None; Raft handles it and the migration is unaffected. \\
    Donor leader       & For each group with \texttt{$G_i$\_START} but no \texttt{$G_i$\_DONE}, the new leader queries the recipient: it commits \texttt{$G_i$\_DONE} if $G_i$ is complete, waits if the data was transferred, and re-issues the RDMA writes if it was only prepared (Algorithm~\ref{alg:donor-recovery}). \\
    Recipient follower & The leader replaces it in the chain with the next live replica (Algorithm~\ref{alg:recipient-node-recovery}); it catches up on the slow path when it recovers. \\
    Recipient leader   & For each group with \texttt{TXN\_START($G_i$)} but no \texttt{TXN\_DONE($G_i$)}, the new leader reports the group's state to the donor leader, which retransmits the data the recipient does not hold durably; data it lacks after \texttt{INDEX\_UPDATE($G_i$)} it fetches from its group (Algorithm~\ref{alg:recipient-leader-recovery}). \\
    Coordinator        & \openq{item 10} \\
    \hline
  \end{tabular}
\end{table}

\subsection{Correctness}
\label{sec:correctness}

\paragraph{No committed write is lost.}
From the start of $G_i$'s migration, writes to $G_i$ are committed only through the recipient's Raft log, and stale-entry detection and tombstones prevent the migrated data from overwriting them.

\paragraph{No read observes partial state.}
Before \texttt{INDEX\_UPDATE($G_i$)} commits, double reads never depend on the recipient alone; after it commits, the durability invariant guarantees that the recipient group holds the data.

\paragraph{Migrations roll forward.}
Every group in flight has a committed start entry without a matching done entry in the donor's or the recipient's log.
A newly elected leader of either group resumes it from that entry (Table~\ref{tab:failures}), and resending data is harmless, so every group eventually reaches \texttt{$G_i$\_DONE}.
\openq{item 10: coordinator failure, and both leaders failing at once.}
